\documentclass{yathesis}

\usepackage[T1]{fontenc}
\usepackage[utf8]{inputenc}
\usepackage{booktabs}
\usepackage{natbib}

\author[]{Jeanne}{Dupont}
\title[]{Tol\'erance aux partitions r\'eseau dans les protocoles de consensus approch\'e}
\academicfield[]{Informatique}
\speciality[]{Syst\`emes distribu\'es}
\date{15}{9}{2026}
\institute[]{Universit\'e de Sciences et Techniques}
\doctoralschool[]{\'Ecole doctorale Informatique, T\'el\'ecommunications et \'Electronique}
\laboratory[]{Laboratoire d'Informatique}{%
  \\
  4 place Jussieu\\
  75005 Paris\\
  France\\
  }
\supervisor[]{Marie}{Martin}
\referee[]{Paul}{Bernard}
\referee[]{Sophie}{Petit}
\committeepresident[]{Luc}{Moreau}
\examiner[]{Claire}{Simon}
\keywords{syst\`emes distribu\'es, consensus, partition r\'eseau}{distributed systems, consensus, network partition}

\begin{document}

\maketitle

\makekeywords

\begin{abstract}
Cette th\`ese \'etudie la d\'egradation des protocoles de consensus
approch\'e sous partition r\'eseau partielle, et propose un protocole
qui \'echange un faible surco\^ut de latence contre un d\'ebit accru.
\end{abstract}
\begin{abstract}
This thesis studies how approximate consensus protocols degrade under
partial network partition, and proposes a protocol that trades a small
latency overhead for improved throughput.
\end{abstract}
\makeabstract

\tableofcontents[depth=chapter,name=Sommaire]
\listoftables
\listoffigures

\mainmatter

\chapter*{Introduction}
\addcontentsline{toc}{chapter}{Introduction}

Les syst\`emes distribu\'es \`a grande \'echelle s'appuient de plus en
plus sur des protocoles de consensus approch\'e pour \'equilibrer
latence et coh\'erence \citep{lamport1998part}. Ce travail \'etudie la
d\'egradation de tels protocoles sous partition r\'eseau.

\chapter{\'Etat de l'art}
\label{ch:etat-art}

\section{Protocoles de consensus classiques}

Les protocoles de consensus classiques comme Paxos
\citep{lamport1998part} et Raft \citep{ongaro2014search} supposent un
r\'eseau totalement connect\'e. Cette th\`ese rel\^ache cette
hypoth\`ese.

\section{Tol\'erance aux partitions}

Les travaux ant\'erieurs sur la tol\'erance aux partitions
sacrifient g\'en\'eralement la disponibilit\'e au profit de la
coh\'erence, suivant le th\'eor\`eme CAP.

\chapter{\'Evaluation}
\label{ch:evaluation}

\section{Protocole exp\'erimental}

Les exp\'eriences ont \'et\'e men\'ees sur un cluster de cinq
n\oe{}uds. Le tableau~\ref{tab:resultats} r\'esume le d\'ebit obtenu
pour des taux de partition croissants.

\begin{table}[htbp]
  \centering
  \caption{D\'ebit (requ\^etes/s) selon le taux de partition simul\'e.}
  \label{tab:resultats}
  \begin{tabular}{@{}lrr@{}}
    \toprule
    Taux de partition & R\'ef\'erence & Propos\'e \\
    \midrule
    0\%  & 12\,400 & 12\,100 \\
    5\%  & 9\,800  & 11\,600 \\
    10\% & 6\,200  & 10\,900 \\
    \bottomrule
  \end{tabular}
\end{table}

\section{Latence sous charge}

La figure~\ref{fig:latence} montre la latence m\'ediane des requ\^etes
lorsque le taux de partition augmente.

\begin{figure}[htbp]
  \centering
  \rule{0.6\textwidth}{4cm}
  \caption{Latence m\'ediane en fonction du taux de partition simul\'e.}
  \label{fig:latence}
\end{figure}

\section{Discussion}

Le protocole propos\'e maintient un meilleur d\'ebit que la
r\'ef\'erence d\`es que le taux de partition d\'epasse 5~\%, au prix
d'un faible surco\^ut de latence constant \citep{ongaro2014search}.


\chapter*{Conclusion}
\addcontentsline{toc}{chapter}{Conclusion}

Le protocole propos\'e maintient un d\'ebit plus \'elev\'e que la
r\'ef\'erence d\`es que le taux de partition d\'epasse 5~\%, au prix
d'un faible surco\^ut de latence constant \citep{ongaro2014search}.

\bibliographystyle{plainnat}
\bibliography{references}

\end{document}
